\begin{table}[ht]
\centering
\begin{threeparttable}
\caption{Returns at stake per filing-year cohort, learning model (millions of 2024 PPP dollars)}\label{tab:levers-learning}
\small
% Generated by em-patent-value code/38_paper_exhibits.py (policy_levers_learning.csv, step 26b); do not edit by hand
\begin{tabular}{lccccc}
\toprule
 & \multicolumn{3}{c}{Fee relief} & Markets for technology & University patenting \\
\cmidrule(lr){2-4}\cmidrule(lr){5-5}\cmidrule(lr){6-6}
 & China & Korea & Brazil & China & China \\
\midrule
Small firms (1--9)          & 0.11 & 0.34 & 0.03 & 797 &  \\
Mid-sized firms (10--99)    & 0.04 & 0.31 & 0.02 & 1{,}194 &  \\
Large firms (100+)          & n.e. & n.e. & n.e. & 0.00 &  \\
Universities                & 0.40 & 0.26 & 0.02 & 4{,}352 & 1{,}009--1{,}460 \\
\bottomrule
\end{tabular}
\begin{tablenotes}\footnotesize
\item \emph{Notes:} As \PaperRef{Table~3}{Table~\ref{tab:levers}}, with the learning model ($d = 0.85$) on 3{,}000 patents per owner group drawn from the estimation cells; owners renew while the expected value of keeping the patent, including the option of later renewals, exceeds the fee, so each lever also changes that option. Markets for technology are calibrated to the same arm's-length survival targets as in \PaperRef{Table~3}{Table~\ref{tab:levers}}, which in this model require far larger changes in returns; the levels are not identified and are shown only to compare the levers. n.e.: not eligible.
\end{tablenotes}
\end{threeparttable}
\end{table}
